% Options for packages loaded elsewhere
\PassOptionsToPackage{unicode}{hyperref}
\PassOptionsToPackage{hyphens}{url}
\documentclass[
]{article}
\usepackage{xcolor}
\usepackage{amsmath,amssymb}
\setcounter{secnumdepth}{-\maxdimen} % remove section numbering
\usepackage{iftex}
\ifPDFTeX
  \usepackage[T1]{fontenc}
  \usepackage[utf8]{inputenc}
  \usepackage{textcomp} % provide euro and other symbols
\else % if luatex or xetex
  \usepackage{unicode-math} % this also loads fontspec
  \defaultfontfeatures{Scale=MatchLowercase}
  \defaultfontfeatures[\rmfamily]{Ligatures=TeX,Scale=1}
\fi
\usepackage{lmodern}
\ifPDFTeX\else
  % xetex/luatex font selection
\fi
% Use upquote if available, for straight quotes in verbatim environments
\IfFileExists{upquote.sty}{\usepackage{upquote}}{}
\IfFileExists{microtype.sty}{% use microtype if available
  \usepackage[]{microtype}
  \UseMicrotypeSet[protrusion]{basicmath} % disable protrusion for tt fonts
}{}
\makeatletter
\@ifundefined{KOMAClassName}{% if non-KOMA class
  \IfFileExists{parskip.sty}{%
    \usepackage{parskip}
  }{% else
    \setlength{\parindent}{0pt}
    \setlength{\parskip}{6pt plus 2pt minus 1pt}}
}{% if KOMA class
  \KOMAoptions{parskip=half}}
\makeatother
\setlength{\emergencystretch}{3em} % prevent overfull lines
\providecommand{\tightlist}{%
  \setlength{\itemsep}{0pt}\setlength{\parskip}{0pt}}
\usepackage{bookmark}
\IfFileExists{xurl.sty}{\usepackage{xurl}}{} % add URL line breaks if available
\urlstyle{same}
\hypersetup{
  hidelinks,
  pdfcreator={LaTeX via pandoc}}

\author{}
\date{}

\begin{document}

\section{NovaCare Final Project
Report}\label{novacare-final-project-report}

\subsection{1. Project Overview}\label{project-overview}

NovaCare is an integrated AI-powered healthcare companion and robotic
assistant designed to empower independence for individuals with physical
or sensory disabilities. The system combines modern web technologies,
machine learning, and physical robotics (SERBot) to create a seamless,
multimodal interaction experience.

The core mission of NovaCare is to bridge communication gaps (via Sign
Language Recognition and Text-to-Speech) and provide intelligent
assistance through Conversational AI and Emotion Detection.

\subsection{2. System Architecture}\label{system-architecture}

The NovaCare repository is organized into a domain-driven,
microservice-based architecture to ensure scalability, isolated
development, and maintainability.

\subsubsection{2.1 Codebase Structure}\label{codebase-structure}

\begin{itemize}
\tightlist
\item
  \textbf{\texttt{apps/}}: End-user interfaces (Next.js Dashboard,
  Flutter Mobile App, React Robot UI).
\item
  \textbf{\texttt{services/}}: Core backend microservices (ASL, LLM,
  Robot HAL, TTS, Optimized Edge Runtime).
\item
  \textbf{\texttt{shared/}}: Common data models, database schemas,
  utility functions, and configurations.
\item
  \textbf{\texttt{infrastructure/}}: Deployment scripts, Docker
  orchestration (\texttt{docker-compose.yml}), and database migrations.
\end{itemize}

\begin{center}\rule{0.5\linewidth}{0.5pt}\end{center}

\subsection{3. American Sign Language (ASL) Recognition
Subsystem}\label{american-sign-language-asl-recognition-subsystem}

A major component of NovaCare is its real-time ASL recognition pipeline,
which allows non-verbal users to communicate naturally with the robot.

\subsubsection{3.1 Model Architecture}\label{model-architecture}

The ASL recognition model utilizes an \textbf{Attention-based Multilayer
Perceptron (MLP)}. - \textbf{Input}: The system uses Google MediaPipe to
extract 21 3D hand landmarks (63 continuous numerical features) per
video frame. - \textbf{Attention Mechanism}: The network treats each of
the 21 landmarks as a ``token'' and applies Multi-Head Self-Attention.
This allows the model to learn the complex spatial relationships between
different fingers (e.g., how the thumb relates to the index finger for
specific letters). - \textbf{Output}: The network predicts across 29
classes (A-Z alphabet, Space, Delete, and an ``Unknown/Idle'' state).

\subsubsection{3.2 Data and Training}\label{data-and-training}

The model was trained on a robust dataset of \textbf{25,479 real-world
samples}. - \textbf{Preprocessing}: Landmarks are normalized relative to
the wrist (landmark 0) to ensure the model is invariant to where the
hand is located in the camera frame. The coordinates are scaled by the
distance to the middle finger (landmark 9) to account for hands being
closer or further from the camera. - \textbf{Augmentation}: During
training, the landmarks undergo synthetic augmentations including random
3D rotations, scaling, translation, and Gaussian noise injection to
prevent overfitting and improve real-world robustness. -
\textbf{Performance}: The final trained model achieved a \textbf{91.59\%
validation accuracy} and \textbf{91.49\% test accuracy} on held-out
data.

\subsubsection{3.3 Inference and Frontend
Integration}\label{inference-and-frontend-integration}

\begin{itemize}
\tightlist
\item
  \textbf{Real-time Pipeline}: The Next.js frontend captures webcam
  frames and sends base64-encoded images to the FastAPI backend.
\item
  \textbf{Temporal Stabilization}: To prevent noisy, flickering
  predictions, the system requires a letter to be continuously predicted
  with high confidence for \textbf{1.5 seconds} before it is
  ``confirmed''.
\item
  \textbf{Accumulator}: Confirmed letters are added to a server-side
  text accumulator. The user can build a full word or sentence and send
  it to the Conversational LLM for a response.
\end{itemize}

\begin{center}\rule{0.5\linewidth}{0.5pt}\end{center}

\subsection{4. Conversational AI and Emotion Detection (LLM)
Subsystem}\label{conversational-ai-and-emotion-detection-llm-subsystem}

The Conversational AI module acts as the cognitive center of NovaCare,
allowing natural language interaction and emotional intelligence.

\subsubsection{4.1 Architecture and APIs}\label{architecture-and-apis}

\begin{itemize}
\tightlist
\item
  \textbf{Framework}: Built on Flask (port 5000), exposing RESTful
  endpoints for chat, history management, and context updates.
\item
  \textbf{LLM Engine}: Integrates with Large Language Models (LLMs) to
  process user prompts, maintain conversation history, and generate
  context-aware responses.
\item
  \textbf{Prompt Engineering}: System prompts are carefully designed to
  instruct the AI to act as a empathetic, concise, and helpful medical
  companion.
\end{itemize}

\subsubsection{4.2 Emotion Detection}\label{emotion-detection}

\begin{itemize}
\tightlist
\item
  \textbf{Background Polling Architecture}: To avoid hardware locking conflicts between backend services and the frontend ASL camera, emotion detection was shifted to a silent frontend background poller. The Next.js frontend uses a custom React hook (\texttt{useBackgroundEmotion}) to capture invisible frames every 5 seconds and evaluates them against a Hugging Face ViT model (\texttt{trpakov/vit-face-expression}).
\item
  \textbf{Live Context Injection}: The frontend dynamically injects the latest detected emotion and confidence score directly into the \texttt{POST\ /api/chat} payload.
\item
  \textbf{Empathy Engine}: The \texttt{MentalHealthOrchestrator} and Conversational AI explicitly ingest this non-verbal data into the LLM's Live Context window (RAG). If a user says "I am fine" but their facial expression reads "sad", the LLM is contextually aware of the contradiction and tailors its empathy and response tone accordingly.
\item
  \textbf{State Emitting}: Based on the conversation and facial analysis, the LLM emits emotional states (e.g., \texttt{happy}, \texttt{concerned}, \texttt{thinking}) which are broadcasted to the Robot UI to physically change the robot's animated facial expressions, bridging the gap between digital AI and physical empathy.
\end{itemize}

\begin{center}\rule{0.5\linewidth}{0.5pt}\end{center}

\subsection{5. Text-to-Speech (TTS)
Subsystem}\label{text-to-speech-tts-subsystem}

To provide a voice for NovaCare, the TTS subsystem converts the LLM's
text responses into natural-sounding audio.

\subsubsection{5.1 Architecture}\label{architecture}

\begin{itemize}
\tightlist
\item
  \textbf{Framework}: A standalone proxy service running on port 8002.
\item
  \textbf{Edge-Optimized}: Designed to be lightweight enough to run on
  edge hardware without requiring massive local deep learning models or
  expensive cloud API calls.
\item
  \textbf{Integration}: Works asynchronously within the Optimized Edge
  Runtime. When the LLM generates a response, it is immediately queued
  to the TTS service, which streams audio to the robot's physical
  speakers.
\end{itemize}

\begin{center}\rule{0.5\linewidth}{0.5pt}\end{center}

\subsection{6. Robot Hardware Abstraction Layer
(HAL)}\label{robot-hardware-abstraction-layer-hal}

The SERBot relies on the Robot HAL to abstract complex motor and
hardware commands into simple, reliable REST APIs.

\subsubsection{6.1 Architecture and
Features}\label{architecture-and-features}

\begin{itemize}
\tightlist
\item
  \textbf{Framework}: Python-based service running on port 9000.
\item
  \textbf{Motor Control}: Exposes endpoints for physical navigation
  (\texttt{move\_forward}, \texttt{turn}, \texttt{stop}), abstracting
  away the low-level GPIO or serial commands required to drive the
  SERBot's motors.
\item
  \textbf{Peripheral Management}: Handles LED controls and physical
  speaker output.
\item
  \textbf{Safety}: Includes built-in safety checks to prevent hardware
  damage, providing a safe testing environment for high-level software.
\end{itemize}

\begin{center}\rule{0.5\linewidth}{0.5pt}\end{center}

\subsection{7. The Optimized Edge
Runtime}\label{the-optimized-edge-runtime}

At the heart of the robot's physical operation is the \textbf{Optimized
Runtime}, an asynchronous orchestrator that ties all microservices
together on the physical robot.

\subsubsection{7.1 Architecture}\label{architecture-1}

\begin{itemize}
\tightlist
\item
  \textbf{Async Orchestrator}: Built entirely around Python's
  \texttt{asyncio}, managing concurrent tasks without blocking the main
  thread.
\item
  \textbf{Task Queues}: Implements priority-based task pipelines for
  audio capture, camera frames, and AI inference.
\item
  \textbf{State Management}: Maintains a thread-safe, unified state of
  the robot's hardware (battery, CPU), current emotion, and conversation
  context.
\item
  \textbf{WebSocket Communication}: Broadcasts real-time state updates
  at 10Hz to the Robot UI.
\item
  \textbf{Non-Destructive Adapters}: Wraps the existing HTTP REST APIs
  (ASL, LLM, TTS, Robot HAL) into the async pipeline, acting as thin,
  fault-tolerant clients that do not modify the original microservices.
\end{itemize}

\begin{center}\rule{0.5\linewidth}{0.5pt}\end{center}

\subsection{8. End-User Applications}\label{end-user-applications}

\subsubsection{8.1 Next.js Web Dashboard}\label{next.js-web-dashboard}

\begin{itemize}
\tightlist
\item
  \textbf{Role}: The primary control center for guardians, medical
  professionals, and patients.
\item
  \textbf{Features}: Real-time ASL recognition interface, vital sign
  monitoring, chat interface, and robot status dashboards.
\end{itemize}

\subsubsection{8.2 Flutter Mobile App}\label{flutter-mobile-app}

\begin{itemize}
\tightlist
\item
  \textbf{Role}: A companion app for on-the-go access.
\item
  \textbf{Features}: Push notifications, remote monitoring, and direct
  communication with the NovaCare system.
\end{itemize}

\subsubsection{8.3 React Robot UI}\label{react-robot-ui}

\begin{itemize}
\tightlist
\item
  \textbf{Role}: The ``face'' of the SERBot, displayed on a physical
  screen mounted on the robot.
\item
  \textbf{Features}: Animated eyes that change color and shape based on
  the current emotion state (e.g., happy blink, thinking rotation). It
  also includes audio level visualizations and connection indicators,
  entirely driven by the WebSocket stream from the Edge Runtime.
\end{itemize}

\begin{center}\rule{0.5\linewidth}{0.5pt}\end{center}

\subsection{9. Deployment and DevOps}\label{deployment-and-devops}

NovaCare is designed to be easily deployable both locally for
development and physically onto the SERBot hardware.

\subsubsection{9.1 Containerization and
Scripts}\label{containerization-and-scripts}

\begin{itemize}
\tightlist
\item
  \textbf{Dockerization}: All microservices are containerized. There are
  separate Docker profiles (\texttt{Dockerfile.serbot},
  \texttt{Dockerfile.laptop}) for edge deployment (optimized for
  Jetson/Raspberry Pi hardware constraints) and heavy laptop-based
  inference.
\item
  \textbf{Unified Scripts}: The project features a unified set of
  PowerShell and Bash scripts (\texttt{start\_all.ps1},
  \texttt{deploy\_serbot.sh}) that handle spinning up the entire
  microservice ecosystem, managing virtual environments, and deploying
  code to the physical robot via SSH/SCP.
\item
  \textbf{Health Monitoring}: Automated scripts
  (\texttt{health\_check.sh}) continuously verify the availability of
  all endpoints and ports.
\end{itemize}

\begin{center}\rule{0.5\linewidth}{0.5pt}\end{center}

\subsection{10. The Team}\label{the-team}

NovaCare was brought to life by a dedicated team of engineers: -
\textbf{Basant Awad} - \textbf{Nadira El-Sirafy} - \textbf{Noureen
Yasser} - \textbf{Muhammad Mustafa} - \textbf{Ramez Asaad}

\begin{center}\rule{0.5\linewidth}{0.5pt}\end{center}

\emph{Generated: June 2026}

\end{document}
